% !TEX root = ../algebra_02.tex
\section{Теорема Коши}

\begin{Theorem}{Коши}{cauchy}
	Пусть $G$ — конечная группа, $p \mid |G|$, $p$ — простое.
	Тогда $\exists\, g \in G\colon \mathrm{ord}\, g = p$.
\end{Theorem}

\begin{proof}
	Рассмотрим множество
	\[
	X = \{(g_1, \ldots, g_p) \in G^p \mid g_1 \cdots g_p = e\}.
	\] 
	Группа $\mathbb{Z}/p\mathbb{Z}$ действует на $X$ циклическим сдвигом:
	\[
	k \cdot (g_1, \ldots, g_p) = (g_{1+k \bmod p},\; g_{2+k \bmod p},\; \ldots,\; g_{p+k \bmod p}).
	\] 

	Заметим, что первые $p-1$ элементов $(g_1, \dots, g_{p-1})$ можно выбрать произвольно из группы $G$, а последний элемент определяется однозначно как $g_p = (g_1 \cdots g_{p-1})^{-1}$. Значит, мощность множества $X$ равна $|X| = |G|^{p-1}$. Так как по условию $p \mid |G|$, то $p$ делит и мощность множества $|X|$.

	Орбиты действия группы $\mathbb{Z}/p\mathbb{Z}$ могут иметь длину, делящую порядок действующей группы, то есть $p$. Так как $p$ — простое число, возможные длины орбит равны либо $1$, либо $p$.
	Пусть $N_1$ — количество орбит длины $1$, а $N_p$ — количество орбит длины $p$.
	Тогда мы можем записать $|X| = N_1 \cdot 1 + N_p \cdot p$.
	Поскольку $p \mid |X|$, из этого уравнения следует, что $p \mid N_1$.

	Орбита состоит из одного элемента тогда и только тогда, когда набор не меняется при циклическом сдвиге, то есть когда $g_1 = g_2 = \dots = g_p = g$. Условие принадлежности такого набора к $X$ означает, что $g^p = e$.
	Один такой элемент мы знаем точно: это набор из нейтральных элементов $(e, e, \dots, e)$ (где $e^p = e$). Следовательно, количество орбит единичной длины $N_1 \ge 1$.
	Так как $N_1$ делится на $p$ и $N_1 \ge 1$, то на самом деле $N_1 \ge p \ge 2$.

	Значит, существует хотя бы один набор $(g, \dots, g)$ с $g \neq e$, для которого $g^p = e$.
	Это означает, что мы нашли неединичный элемент $g \in G$, порядок которого равен $p$. Теорема доказана.
\end{proof}
